\subsection{Decision Tree and Implementation}

The game maps a single industry pathway (marketing) from 2025 to 2035. After a shared opening that introduces the player to their team and to Alex, the player's manager, the game presents five yearly decisions from 2025 through 2029. Each decision is staged as a brief news beat, a short context scene, and a multiple-choice prompt, and the same five decisions appear on every playthrough; what differs is the option the player selects. Each option silently increments one or two of four hidden path counters corresponding to the four plausible futures described above. Two visible meters --- the player's standing with their manager and the team's read on the AI agenda --- also shift in response to choices, but do not by themselves determine the ending; they are surfaced so that the social and political texture of each choice is felt rather than abstracted.

The ending is resolved at the 2030 ballot on the AI Job Preservation and Work Sharing Act. A \textsc{for} vote always passes the bill, an \textsc{against} vote always defeats it, and an \textsc{abstain} vote passes only if the player's augmentation and pro-worker counters together exceed the shrinkage and erosion counters. The winning side of the ballot then selects the ending by comparing the two counters on that side. A small reproducible perturbation (approximately 15\% of possible playthroughs) occasionally flips the ending to its sibling on the same side, so that runs landing near the boundary can experience either outcome. This perturbation is computed from a deterministic seeded hash of the player's state rather than a true random draw, so the same playthrough always produces the same ending; page refreshes, hot reloads, and the resume feature cannot change the outcome mid-run. From 2031 to 2035 each ending carries its own narrative arc, several of which include additional sub-decisions, and the game closes on a results screen that visualizes the player's full path through the decision tree. From that screen, players can also open a short policies appendix surfacing five concrete real-world levers --- work-sharing legislation, an AI dividend, human-in-the-loop standards, portable retraining accounts, and workplace metric transparency --- each drawn from the labor and AI policy research reviewed above, so that the player's in-game outcome connects directly to the civic and policy questions they could engage with outside the game.

The game is built as a browser-based interactive experience with an illustrative, pixel-art visual style rather than a photorealistic one, drawing on the tradition of narrative games that use stylized, atmospheric artwork to carry serious subject matter. Our reference points include games which use a deliberately constrained visual style to create moral weight through bureaucratic decisions, and web-based interactive experiences like the \textit{Financial Times} Climate Game and Nicky Case's \textit{The Evolution of Trust}. We aimed for a visual register that signals ``this is a game'' immediately upon loading, priming players toward the role-play engagement described above, while remaining achievable within our team's capacity and timeline.

The application is built in Next.js and React with TypeScript. Rather than adopting an external interactive-fiction engine, we encode the entire decision graph as a single typed data structure listing every year, every choice, and the effects each choice has on the hidden counters and visible meters; scene components, the in-game futures diagram, and the analytics layer all read from this same structure, so gameplay and visualization cannot drift apart. Game state is managed by a pure reducer, persisted to the browser's local storage with a versioned save format that supports resuming an in-progress run, and surfaced to scenes through React context. Anonymous per-session analytics are captured through PostHog using a small set of gameplay events (\texttt{game\_started}, \texttt{decision\_made}, \texttt{policy\_voted}, \texttt{ending\_reached}, and a few others) along with the choice made at each branch point, allowing us to study aggregate behavior without collecting any personally identifying information. The results-screen tree visualization is a custom SVG view that reads from the same decision-graph data structure used during play. The game is deployed as a static site and is shareable via a single URL, which supports our goal of broad public accessibility.
